\documentclass[12pt, a4paper]{amsart}

% ── Packages ──────────────────────────────────────────────────────────────────
\usepackage{amsmath, amssymb, amsthm}
\usepackage{mathtools}
\usepackage{geometry}
\usepackage{hyperref}
\usepackage{microtype}
\usepackage{enumitem}
\usepackage{array}
\usepackage{booktabs}
\usepackage{xcolor}

\geometry{
  top    = 2.8cm,
  bottom = 3.0cm,
  left   = 3.0cm,
  right  = 3.0cm
}

\hypersetup{
  colorlinks = true,
  linkcolor  = blue!60!black,
  citecolor  = blue!60!black,
  urlcolor   = blue!60!black
}

% ── Theorem environments ───────────────────────────────────────────────────────
\theoremstyle{plain}
\newtheorem{theorem}{Theorem}[section]
\newtheorem{proposition}[theorem]{Proposition}
\newtheorem{lemma}[theorem]{Lemma}
\newtheorem{corollary}[theorem]{Corollary}

\theoremstyle{definition}
\newtheorem{definition}[theorem]{Definition}
\newtheorem{example}[theorem]{Example}
\newtheorem{remark}[theorem]{Remark}

\theoremstyle{remark}
\newtheorem{conjecture}[theorem]{Conjecture}
\newtheorem{question}[theorem]{Open Question}

% ── Macros ─────────────────────────────────────────────────────────────────────
\newcommand{\N}{\mathbb{N}}
\newcommand{\Z}{\mathbb{Z}}
\newcommand{\Q}{\mathbb{Q}}
\newcommand{\R}{\mathbb{R}}
\newcommand{\C}{\mathbb{C}}
\newcommand{\Zs}{\Z[\tfrac{1}{6}]}
\newcommand{\Aff}{\operatorname{Aff}}
\newcommand{\val}{\nu}
\newcommand{\Stab}{\operatorname{Stab}}
\newcommand{\wL}{\mathsf{L}}
\newcommand{\wR}{\mathsf{R}}

% ── Title block ────────────────────────────────────────────────────────────────
\title[Collatz as Orbit Closure in $\Aff(\Zs)$]{%
  The Collatz Conjecture as an Orbit-Closure Problem\\
  in the Affine Group $\Aff\!\left(\Z\!\left[\tfrac{1}{6}\right]\right)$
}

\author{Research Notes}
\date{\today}

\begin{document}

\begin{abstract}
We reformulate the Collatz conjecture as an orbit-closure statement in
the affine group $\Aff(\Zs)$.  Every word $w$ in the two generators
$\wL\colon x\mapsto x/2$ and $\wR\colon x\mapsto 3x+1$ acts as an
affine map $f_w(x)=ax+b$ with slope $a=3^r/2^\ell$ and intercept
$b\in\Zs$.  We prove that the Collatz semigroup
$\mathcal{S}=\langle\wL,\wR\rangle$ embeds faithfully into $\Aff(\Zs)$
(Theorem~\ref{thm:faithful}), characterise the word-equality relation
(Theorem~\ref{thm:word-equality}), and show the conjecture is
equivalent to the backward orbit of $1$ under $\mathcal{S}$ covering
all of $\N$ (Theorem~\ref{thm:equiv}).  We establish that $\Aff(\Zs)$
is solvable and hence amenable (Proposition~\ref{prop:amenable}),
identifying this as a structural advantage over Thompson-group
approaches where amenability of the ambient group is itself an open
problem.  The residual gap between density-$1$ coverage and full
coverage is identified precisely, and the conjecture is reduced to a
concrete subset-sum reachability question over $\Zs$
(Question~\ref{q:reachability}).
\end{abstract}

\maketitle
\tableofcontents

%% ═══════════════════════════════════════════════════════════════════════════
\section{Introduction}
%% ═══════════════════════════════════════════════════════════════════════════

The Collatz conjecture asserts that iterating the map
\[
  T(n) \;=\; \begin{cases}
    n/2   & \text{if } n \text{ is even,} \\
    3n+1  & \text{if } n \text{ is odd,}
  \end{cases}
\]
from any positive integer $n$ eventually reaches $1$.
Despite its elementary statement the problem has resisted all attempts
at proof since Collatz posed it informally around 1937.  The strongest
partial result is due to Tao~\cite{Tao2019}, who showed that the
Collatz orbit of almost every positive integer (in the logarithmic
density sense) reaches arbitrarily small values.

This paper develops a new algebraic framework centred on the
\emph{affine group} $\Aff(\Zs)$ of transformations
$x\mapsto ax+b$ over the ring $\Zs=\Z[\frac{1}{6}]$.  The even and
odd branches of $T$ are each affine maps, every Collatz orbit prefix
composes to a single affine map, and the Collatz conjecture becomes a
question about orbit coverage under a semigroup action---a setting with
clear algebraic structure.

\subsection*{Main results}

\begin{itemize}
  \item \textbf{Faithful embedding} (Section~\ref{sec:embedding}).
        The Collatz semigroup $\mathcal{S}=\langle\wL,\wR\rangle$ embeds
        faithfully into $\Aff(\Zs)$.  Two words are equal as semigroup
        elements if and only if they have the same number of each
        generator \emph{and} the same intercept.  All
        $\binom{r+\ell}{r}$ orderings of a word with $r$ copies of
        $\wR$ and $\ell$ copies of $\wL$ produce distinct affine maps.
        The word problem is decidable.

  \item \textbf{Orbit reformulation} (Section~\ref{sec:orbit}).
        The conjecture is equivalent to: for every $n\in\N$, there
        exists $f\in\mathcal{S}$ with $f(n)=1$.

  \item \textbf{Amenability} (Section~\ref{sec:amenable}).
        $\Aff(\Zs)$ is solvable of derived length $2$, hence amenable
        by Day's theorem~\cite{Day1957}.  This gives density-$1$
        coverage for free, recovering Tao's result in this language.

  \item \textbf{The residual gap} (Section~\ref{sec:gap}).
        Amenability alone does not close the gap between density-$1$
        and full coverage.  The precise obstruction is reduced to a
        subset-sum reachability problem over $\Zs$.
\end{itemize}

\subsection*{Comparison with Thompson-group approaches}

Several authors have explored embeddings of the Collatz dynamics into
Thompson's groups $F$, $T$, and $V$
\cite{Cleary2000,BrunnerSidki1999}.  A key difficulty is that
amenability of Thompson's group $F$ is a major open problem
\cite{BelkBrown2005,Haagerup2016}.  Any proof relying on amenability
of the ambient group therefore cannot use $F$ without first resolving
that problem.  The affine group $\Aff(\Zs)$ is \emph{provably} amenable,
making the present framework structurally simpler.

The related group $G_{(2,3)}$ of piecewise-linear homeomorphisms with
breakpoints in $\Zs$ and slopes in $\langle 2,3\rangle$ is discussed in
Section~\ref{sec:further} as a direction for more rigid embedding.

%% ═══════════════════════════════════════════════════════════════════════════
\section{Setup and notation}
\label{sec:setup}
%% ═══════════════════════════════════════════════════════════════════════════

\begin{definition}[The ring $\Zs$]
  Write $\Zs = \Z[\tfrac{1}{6}]$ for the smallest subring of $\Q$
  containing $\Z$ and closed under division by $2$ and by $3$:
  \[
    \Zs \;=\;
    \Bigl\{ \frac{p}{2^a 3^b} : p\in\Z,\; a,b\geq 0 \Bigr\}.
  \]
  The units of $\Zs$ are $\{\pm 2^s 3^t : s,t\in\Z\}$.
\end{definition}

\begin{definition}[Affine group]
  The \emph{orientation-preserving affine group} is
  \[
    \Aff^+(\Zs) \;=\;
    \bigl\{ x\mapsto ax+b : a = 2^s 3^t > 0,\; s,t\in\Z,\; b\in\Zs \bigr\},
  \]
  with composition $(a_1,b_1)\circ(a_2,b_2) = (a_1 a_2,\, a_1 b_2 + b_1)$,
  identity $(1,0)$, and inverse $(a,b)^{-1}=(a^{-1},-b/a)$.
  We write elements as pairs $(a,b)$ and the corresponding map as $f_{(a,b)}$.
\end{definition}

\begin{definition}[Generators]
  \label{def:generators}
  Define two elements of $\Aff^+(\Zs)$:
  \begin{align*}
    \wL &= \bigl(\tfrac{1}{2},\, 0\bigr), &
    f_\wL(x) &= x/2 \quad \text{(even branch of } T\text{),} \\
    \wR &= (3,\, 1), &
    f_\wR(x) &= 3x+1 \quad \text{(odd branch of } T\text{).}
  \end{align*}
  Their group inverses are
  $\wL^{-1}=(2,0)$ (that is, $x\mapsto 2x$) and
  $\wR^{-1}=(\frac{1}{3},-\frac{1}{3})$ (that is, $x\mapsto (x-1)/3$).
\end{definition}

\begin{definition}[Collatz semigroup and words]
  \label{def:semigroup}
  The \emph{Collatz semigroup} is
  $\mathcal{S} = \langle\wL,\wR\rangle \subset \Aff^+(\Zs)$.
  A \emph{word} is a finite string $w=c_1 c_2\cdots c_n$ with
  $c_i\in\{\wL,\wR\}$; its \emph{type} is $(r,\ell)$ where $r$ counts
  $\wR$'s and $\ell$ counts $\wL$'s.  Composition is left-to-right:
  $f_w = f_{c_n}\circ\cdots\circ f_{c_1}$, so $f_w(x)$ is obtained
  by applying $f_{c_1}$ first.
\end{definition}

\begin{remark}[Canonical word of an orbit]
  \label{rem:canonical}
  The Collatz sequence $n, T(n), T^2(n), \ldots, 1$ determines a
  canonical word $w(n)$: write $\wL$ for each even step and $\wR$ for
  each odd step.  Then $f_{w(n)}(n)=1$ if and only if the sequence
  terminates at $1$.
\end{remark}

%% ═══════════════════════════════════════════════════════════════════════════
\section{The affine embedding}
\label{sec:embedding}
%% ═══════════════════════════════════════════════════════════════════════════

\subsection{Every word is an affine map}

\begin{lemma}[Affine structure]
  \label{lem:affine}
  For any word $w = c_1\cdots c_n$ of type $(r,\ell)$, the element
  $f_w\in\mathcal{S}$ has the form $f_w(x) = a_w x + b_w$ where
  \begin{align}
    a_w &= \frac{3^r}{2^\ell}, \label{eq:slope} \\
    b_w &= \sum_{\substack{i=1\\c_i=\wR}}^{n}
           \frac{3^{r_i}}{2^{\ell_i}}, \label{eq:intercept}
  \end{align}
  and $r_i = \#\{j < i : c_j = \wR\}$,
  $\ell_i = \#\{j < i : c_j = \wL\}$ denote the number of $\wR$'s
  (resp.\ $\wL$'s) strictly before position $i$.
\end{lemma}

\begin{proof}
  By induction on $n$.  The base cases
  $f_\wL(x)=\frac{1}{2}x+0$ and $f_\wR(x)=3x+1$ match
  \eqref{eq:slope}--\eqref{eq:intercept}.

  Suppose $f_v(x) = ax+b$ for a word $v$ of type $(r',\ell')$.

  \noindent\emph{Appending $\wL$:}
  $f_{v\wL}(x) = f_v(x/2) = \frac{a}{2}x + \frac{b}{2}$.
  The new slope is $a/2 = 3^{r'}/2^{\ell'+1}$, consistent with
  \eqref{eq:slope} for type $(r',\ell'+1)$.  Appending $\wL$ adds no
  new term to \eqref{eq:intercept}, consistent with the sum being over
  $\wR$-positions only.

  \noindent\emph{Appending $\wR$:}
  $f_{v\wR}(x) = f_v(3x+1) = 3ax + (a+b)$.
  The new slope is $3a = 3^{r'+1}/2^{\ell'}$, consistent with
  \eqref{eq:slope} for type $(r'+1,\ell')$.  The new intercept is
  $a+b = 3^{r'}/2^{\ell'} + b$, which is $b$ plus the new term
  $3^{r'}/2^{\ell'}$ contributed by the appended $\wR$ at position
  $n+1$ (which has $r' = r'$ preceding $\wR$'s and $\ell'$ preceding
  $\wL$'s), consistent with \eqref{eq:intercept}.
\end{proof}

\subsection{Injectivity}

\begin{theorem}[Faithful embedding]
  \label{thm:faithful}
  The map $\mathcal{S}\to\Aff^+(\Zs)$, $w\mapsto (a_w,b_w)$, is
  injective.  The Collatz semigroup is isomorphic to its image in
  $\Aff^+(\Zs)$.
\end{theorem}

\begin{proof}
  Suppose $f_{w_1}=f_{w_2}$, i.e.\ $a_{w_1}=a_{w_2}$ and
  $b_{w_1}=b_{w_2}$.  By \eqref{eq:slope} and unique factorisation in
  $\Z$, equality of slopes $3^{r_1}/2^{\ell_1}=3^{r_2}/2^{\ell_2}$
  forces $r_1=r_2=r$ and $\ell_1=\ell_2=\ell$.

  Now suppose $w_1\neq w_2$ are two distinct words of the same type
  $(r,\ell)$.  Let $k$ be the first position at which they differ.
  Without loss of generality, $c_k^{(1)}=\wR$ and $c_k^{(2)}=\wL$.
  Both words agree on positions $1,\ldots,k-1$, so the partial counts
  $(r_{k}^{(1)}, \ell_{k}^{(1)}) = (r_{k}^{(2)}, \ell_{k}^{(2)})$
  are identical; call them $(r_0,\ell_0)$.  The term
  $3^{r_0}/2^{\ell_0}\in\Zs$ appears in $b_{w_1}$ (contributed by
  position $k$) but not yet in $b_{w_2}$.  After position $k$, the
  remaining letters of $w_1$ and $w_2$ contribute terms with
  denominators divisible by a power of $2$ at least $\ell_0+1$ (since
  $w_2$ has an $\wL$ at position $k$, incrementing $\ell$), while
  $3^{r_0}/2^{\ell_0}$ has denominator exactly $2^{\ell_0}$.  Since
  elements of $\Zs$ with different denominators are distinct, and since
  the remaining contributions to $b_{w_2}$ cannot cancel the term
  $3^{r_0}/2^{\ell_0}$ present only in $b_{w_1}$, we conclude
  $b_{w_1}\neq b_{w_2}$.
\end{proof}

\subsection{Word equality and the word problem}

\begin{theorem}[Word equality]
  \label{thm:word-equality}
  Two words $w_1, w_2\in\{\wL,\wR\}^*$ represent the same element of
  $\mathcal{S}$ if and only if:
  \begin{enumerate}[label=\textup{(\roman*)}]
    \item $w_1$ and $w_2$ have the same type $(r,\ell)$, and
    \item $b_{w_1} = b_{w_2}$.
  \end{enumerate}
  All $\binom{r+\ell}{r}$ orderings of a type-$(r,\ell)$ word produce
  pairwise distinct affine maps.  The word problem for $\mathcal{S}$ is
  decidable in time $O(|w|)$.
\end{theorem}

\begin{proof}
  Conditions (i)--(ii) are necessary by Theorem~\ref{thm:faithful}.
  They are sufficient by definition of $f_w$.  Distinctness of all
  $\binom{r+\ell}{r}$ orderings follows from the injectivity argument:
  any two distinct orderings of the same $(r,\ell)$-type words differ
  at some first position, and the argument in the proof of
  Theorem~\ref{thm:faithful} shows their intercepts differ.
  Decidability: given words $w_1, w_2$, compute $(a_{w_i},b_{w_i})$
  using Lemma~\ref{lem:affine} in $O(|w_i|)$ time, then compare.
\end{proof}

\begin{example}[First non-trivial relation]
  The words $\mathsf{LRRLRLL}$ and $\mathsf{RRLLLLR}$ both have type
  $(r,\ell)=(3,4)$.  By Lemma~\ref{lem:affine}:
  \[
    f_{\mathsf{LRRLRLL}}(x) = f_{\mathsf{RRLLLLR}}(x)
    = \frac{27}{16}\,x + \frac{7}{4}.
  \]
  This is the shortest non-trivial relation in the semigroup.
  Both words produce the same intercept $b=\frac{7}{4}$, which can be
  verified by computing the contribution of each $\wR$-position.
\end{example}

\begin{corollary}[Group structure]
  \label{cor:structure}
  The semigroup $\mathcal{S}$ satisfies none of the standard relations
  that would identify it with a classical group:
  \[
    f_{\wL\wR}(x) = \tfrac{3}{2}x+1, \quad
    f_{\wR\wL}(x) = \tfrac{3}{2}x+\tfrac{1}{2},
  \]
  so $\wL\wR\neq\wR\wL$ (not commutative).  Similarly,
  $\wL^2\wR\neq\wR\wL$ (no Thompson-$F$ relation), and
  $\wL\wR\wL^{-1}\neq\wR^2$ (no standard conjugation relation).
\end{corollary}

%% ═══════════════════════════════════════════════════════════════════════════
\section{Orbit closure reformulation}
\label{sec:orbit}
%% ═══════════════════════════════════════════════════════════════════════════

\subsection{The fiber map and backward orbit}

\begin{definition}[Fiber at $1$]
  For $f=(a,b)\in\mathcal{S}$, the \emph{fiber at $1$} is the unique
  solution to $f(x)=1$:
  \[
    \phi(f) \;:=\; f^{-1}(1) \;=\; \frac{1-b}{a} \;\in\;\Zs.
  \]
  The \emph{backward orbit of $1$} under $\mathcal{S}$ is
  $\mathcal{O}^-(1) := \{\phi(f) : f\in\mathcal{S}\} \subset \Zs$.
\end{definition}

\begin{definition}[Intercept sets]
  \label{def:intercepts}
  For non-negative integers $r,\ell$, let
  \[
    \mathcal{B}(r,\ell) \;:=\;
    \bigl\{b_w : w\in\{\wL,\wR\}^*,\; w\text{ has type }(r,\ell)\bigr\}
    \;\subset\;\Zs.
  \]
\end{definition}

\subsection{The equivalence theorem}

\begin{theorem}[Four equivalent formulations]
  \label{thm:equiv}
  The Collatz conjecture is equivalent to each of the following:
  \begin{enumerate}[label=\textup{(\Alph*)}]
    \item\label{it:A}
      \emph{Orbit coverage}: For every $n\in\N$, there exists a word
      $w\in\{\wL,\wR\}^*$ such that $f_w(n)=1$.
    \item\label{it:B}
      \emph{Backward orbit}: $\mathcal{O}^-(1)\supseteq\N$.
    \item\label{it:C}
      \emph{Reachability}: For every $n\in\N$, there exist $r,\ell\geq 0$
      such that $1-3^r n/2^\ell \in \mathcal{B}(r,\ell)$.
    \item\label{it:D}
      \emph{Tree coverage}: The inverse tree $\mathcal{T}$ rooted at
      $1$, with branching rules
      \[
        n\mapsto 2n \quad\text{and}\quad
        n\mapsto\tfrac{n-1}{3}
        \;\bigl(\text{when }n\equiv 1\!\pmod{3}\text{ and }
        \tfrac{n-1}{3}\text{ is odd}\bigr),
      \]
      contains every positive integer as a node.
  \end{enumerate}
\end{theorem}

\begin{proof}
  \textbf{Collatz}$\Leftrightarrow$\textbf{\ref{it:A}}:
  If the Collatz sequence for $n$ terminates at $1$, the word $w(n)$
  of Remark~\ref{rem:canonical} satisfies $f_{w(n)}(n)=1$ by
  construction, since each step $c_i$ applies either the even or odd
  branch of $T$ and the composition carries $n$ to $1$.
  Conversely, if $f_w(n)=1$ for some $w=c_1\cdots c_k$, then the
  sequence $n, f_{c_1}(n), f_{c_1 c_2}(n), \ldots, f_w(n)=1$ is a
  valid Collatz trajectory from $n$ to $1$.

  \textbf{\ref{it:A}}$\Leftrightarrow$\textbf{\ref{it:B}}:
  $f_w(n)=1$ if and only if $n=\phi(f_w)\in\mathcal{O}^-(1)$.

  \textbf{\ref{it:B}}$\Leftrightarrow$\textbf{\ref{it:C}}:
  For a type-$(r,\ell)$ word, $\phi(f_w)=n$ rewrites as
  $1-b_w = n\cdot(3^r/2^\ell)$, i.e.\ $b_w = 1-3^r n/2^\ell$.
  Since $b_w\in\mathcal{B}(r,\ell)$ by definition, this is the
  condition in~\ref{it:C}.

  \textbf{\ref{it:B}}$\Leftrightarrow$\textbf{\ref{it:D}}:
  The tree $\mathcal{T}$ is the Cayley graph of
  $\mathcal{O}^-(1)\cap\N$ under the inverse generators
  $\phi(\wL^{-1})(n) = 2n$ and $\phi(\wR^{-1})(n) = (n-1)/3$
  (defined when $n\equiv 1\pmod{3}$ and $(n-1)/3$ is odd, so that
  $T((n-1)/3)=n$ via the odd branch).  Coverage of $\N$ is
  equivalent to $\mathcal{O}^-(1)\supseteq\N$.
\end{proof}

\begin{example}[Verification for $n=7$]
  The Collatz sequence $7\to22\to11\to34\to17\to52\to26\to13
  \to40\to20\to10\to5\to16\to8\to4\to2\to1$ has type $(r,\ell)=(5,11)$.
  By \eqref{eq:slope}--\eqref{eq:intercept},
  \[
    a_{w(7)} = \frac{3^5}{2^{11}} = \frac{243}{2048},
    \qquad
    b_{w(7)} = \frac{347}{2048},
  \]
  and $\frac{243}{2048}\cdot 7 + \frac{347}{2048}
       = \frac{1701+347}{2048} = \frac{2048}{2048} = 1$. \qed
\end{example}

\subsection{Structure of $\mathcal{B}(r,\ell)$}

\begin{proposition}[Intercept sets]
  \label{prop:intercepts}
  For all $r,\ell\geq 0$:
  \begin{enumerate}[label=\textup{(\roman*)}]
    \item $|\mathcal{B}(r,\ell)| = \binom{r+\ell}{r}$: all orderings
          of a type-$(r,\ell)$ word produce distinct intercepts.
    \item Every element of $\mathcal{B}(r,\ell)$ has denominator
          dividing $2^\ell$.
    \item $\mathcal{B}(r,0) = \{(3^r-1)/2\}$ (unique intercept when
          all $\wR$'s precede all $\wL$'s, giving the geometric sum
          $1+3+\cdots+3^{r-1}=(3^r-1)/2$).
    \item $\mathcal{B}(0,\ell) = \{0\}$.
  \end{enumerate}
\end{proposition}

\begin{proof}
  Part (i) is Theorem~\ref{thm:word-equality}.  Part (ii) follows
  from \eqref{eq:intercept}: each term has denominator $2^{\ell_i}
  \leq 2^\ell$.  Parts (iii)--(iv) are direct computations.
\end{proof}

\begin{table}[ht]
  \centering
  \caption{Intercept sets $\mathcal{B}(r,\ell)$ for small $(r,\ell)$.}
  \medskip
  \begin{tabular}{c c c l}
    \toprule
    $r$ & $\ell$ & $\bigl|\mathcal{B}(r,\ell)\bigr|=\binom{r+\ell}{r}$
        & Elements of $\mathcal{B}(r,\ell)$ \\
    \midrule
    0 & 2 & 1 & $\{0\}$ \\
    1 & 1 & 2 & $\bigl\{\tfrac{1}{2},\; 1\bigr\}$ \\
    1 & 2 & 3 & $\bigl\{\tfrac{1}{4},\; \tfrac{1}{2},\; 1\bigr\}$ \\
    1 & 3 & 4 &
      $\bigl\{\tfrac{1}{8},\;\tfrac{1}{4},\;\tfrac{1}{2},\;1\bigr\}$ \\
    2 & 2 & 6 &
      $\bigl\{1,\;\tfrac{5}{4},\;\tfrac{7}{4},\;2,\;\tfrac{5}{2},\;4\bigr\}$ \\
    2 & 3 & 10 &
      $\bigl\{\tfrac{1}{2},\;\tfrac{5}{8},\;\tfrac{7}{8},\;1,\;\tfrac{5}{4},\;
               \tfrac{11}{8},\;\tfrac{7}{4},\;2,\;\tfrac{5}{2},\;4\bigr\}$ \\
    \bottomrule
  \end{tabular}
\end{table}

%% ═══════════════════════════════════════════════════════════════════════════
\section{Amenability of $\Aff^+(\Zs)$}
\label{sec:amenable}
%% ═══════════════════════════════════════════════════════════════════════════

\begin{proposition}[Solvability and amenability]
  \label{prop:amenable}
  The group $\Aff^+(\Zs)$ is solvable of derived length $2$, and
  therefore amenable.
\end{proposition}

\begin{proof}
  Let $T = \{(1,b) : b\in\Zs\}$ be the normal subgroup of translations.
  The quotient $\Aff^+(\Zs)/T$ is isomorphic to the multiplicative
  group $\{2^s 3^t : s,t\in\Z\}\cong\Z^2$, which is abelian.  Hence
  $[\Aff^+(\Zs),\Aff^+(\Zs)]\subseteq T$.
  Since $T\cong(\Zs,+)$ is abelian, $[T,T]=\{0\}$.
  The derived series
  $\Aff^+(\Zs)\supseteq T\supseteq\{e\}$ witnesses solvability of
  length at most $2$.
  Amenability follows from Day's theorem~\cite{Day1957}: abelian groups
  are amenable, and the class of amenable groups is closed under
  extensions.
\end{proof}

\begin{remark}[Thompson's $F$ contrast]
  Thompson's group $F$ is the group of piecewise-linear homeomorphisms
  of $[0,1]$ with finitely many breakpoints in $\Z[\frac{1}{2}]$ and
  slopes in $\{2^k\}$.  Whether $F$ is amenable is open;
  Haagerup--Olesen~\cite{Haagerup2016} showed $F$ is not inner amenable.
  The group $\Aff^+(\Zs)$ is provably amenable, so proofs requiring
  amenability of the ambient group are feasible here but not in the
  Thompson setting.
\end{remark}

\subsection{Negative drift and density-1 coverage}

\begin{proposition}[Expected log-decrease per Syracuse step]
  \label{prop:drift}
  Let $n$ be odd and let $k=\nu_2(3n+1)$.  The composed step
  $\wR\wL^k$ (odd branch followed by $k$ halvings) multiplies $n$ by
  $3/2^k$.  Taking logarithms, the expected change per Syracuse step is
  \[
    \mathbf{E}[\log(3/2^k)]
    = \log 3 - (\log 2)\,\mathbf{E}[k].
  \]
  Under the heuristic that $k=\nu_2(3n+1)$ is geometrically distributed
  with $P(k=j)=1/2^j$, $\mathbf{E}[k]=2$, and
  \[
    \mathbf{E}[\log(3/2^k)] = \log 3 - 2\log 2 = \log(3/4) \approx -0.288.
  \]
\end{proposition}

\begin{remark}
  Combining the negative drift of Proposition~\ref{prop:drift} with the
  ergodic theorem for amenable groups
  (Lindenstrauss~\cite{Lindenstrauss2001}) gives logarithmic density $1$
  for $\mathcal{O}^-(1)$ in $\N$, recovering the quantitative content of
  Tao's theorem~\cite{Tao2019} in this algebraic language.
\end{remark}

%% ═══════════════════════════════════════════════════════════════════════════
\section{The amenability gap and the residual problem}
\label{sec:gap}
%% ═══════════════════════════════════════════════════════════════════════════

\subsection{What amenability gives and does not give}

Amenability of $\Aff^+(\Zs)$ and negative drift together imply:
\[
  \frac{1}{\log N}\sum_{\substack{n\leq N \\ n\notin\mathcal{O}^-(1)}}
  \frac{1}{n} \;\longrightarrow\; 0.
\]
This is a density statement.  It does \emph{not} imply full coverage,
because a set of logarithmic density $1$ may still omit finitely (or
even infinitely many sparse) integers.  We call this the
\emph{amenability gap}.

\subsection{The open problem}

\begin{question}[Subset-sum reachability]
  \label{q:reachability}
  For every $n\in\N$, do there exist $r,\ell\geq 0$ such that
  \[
    1 - \frac{3^r n}{2^\ell} \;\in\; \mathcal{B}(r,\ell)\;?
  \]
\end{question}

An affirmative answer to Question~\ref{q:reachability} is equivalent
to the Collatz conjecture (Theorem~\ref{thm:equiv}\ref{it:C}).

\begin{remark}[Feasibility of the target]
  By Euler's theorem (since $\gcd(3,2^\ell)=1$), for each $n$ and each
  $\ell$ there exists $r$ with $3^r n \equiv 1\pmod{2^\ell}$, making
  $1-3^r n/2^\ell\in\Zs$.  So the condition is non-vacuous: for every
  $n$ and $\ell$ there is at least one $(r,\ell)$ for which the target
  lands in $\Zs$.  The question is whether it lands in the smaller set
  $\mathcal{B}(r,\ell)$.
\end{remark}

\begin{remark}[Density of targets]
  Since $\log 3/\log 2$ is irrational, the set
  $\{3^r/2^\ell : r,\ell\geq 0\}$ is dense in $\R_{>0}$.  For fixed
  $n$, the targets $1-3^r n/2^\ell$ are therefore dense in $\R$ as
  $(r,\ell)$ varies.  Combined with the fact that
  $|\mathcal{B}(r,\ell)|=\binom{r+\ell}{r}\to\infty$, this makes the
  reachability condition plausible but does not prove it.
\end{remark}

%% ═══════════════════════════════════════════════════════════════════════════
\section{The $\operatorname{Im}(\tau)$ dynamics and the polylogarithm}
\label{sec:tau}
%% ═══════════════════════════════════════════════════════════════════════════

A parallel line of investigation, predating the affine embedding,
attaches a complex parameter $\tau(n)$ to each iterate and studies its
dynamics.  The results here are self-contained and complement the
affine picture.

\begin{definition}[$\tau$-parameter]
  For $n\in\N$ let $a_0=\nu_2(n)$ and $a_1=\nu_3(n)$ (the $2$- and
  $3$-adic valuations).  Define
  \[
    \tau(n) \;=\; \frac{a_1 + i\,\nu_2(3n+1)}{a_0+1} \;\in\; \C.
  \]
\end{definition}

\begin{proposition}[Imaginary axis dominance]
  For almost all iterates of the Collatz map, $a_1=0$ and hence
  $\operatorname{Re}(\tau)=0$.  The sequence $\tau(n), \tau(T(n)),\ldots$
  lives almost entirely on the positive imaginary axis, and the
  dynamics reduce to the real sequence
  $y_k := \operatorname{Im}(\tau(T^k(n))) \in \R_{>0}$.
\end{proposition}

\begin{proposition}[Rational transition rules]
  \label{prop:transition}
  The sequence $(y_k)$ obeys:
  \begin{align*}
    \text{Even step } (n\to n/2): &\quad
      y_{k+1} = y_k \cdot \frac{a_0^2-1}{a_0^2}
      = y_k \cdot \bigl(1 - a_0^{-2}\bigr), \\
    \text{Odd step } (n\to 3n+1): &\quad
      y_{k+1} = y_k \cdot \frac{1}{\nu_2(3n+1)+1}.
  \end{align*}
  Every sequence terminates at $y=\tfrac{1}{2}$ (corresponding to $n=2$).
\end{proposition}

\begin{proposition}[$\pm S$ symmetry and the polylogarithm]
  \label{prop:symmetry}
  Define $S = \sum_{k=1}^\infty k(\log k)/2^k \approx 0.5078$.  Then
  \[
    \mathbf{E}[\Delta\log y \mid \text{odd step}] = +S, \qquad
    \mathbf{E}[\Delta\log y \mid \text{even step}] = -S.
  \]
  The constant $S$ equals the negative of the derivative of the
  polylogarithm at $s=0$:
  \[
    S \;=\; -\frac{d}{ds}\operatorname{Li}_s\!\bigl(\tfrac{1}{2}\bigr)
            \bigg|_{s=0}.
  \]
\end{proposition}

\begin{proof}[Proof sketch]
  Under the heuristic that $\nu_2(3n+1)$ is geometrically distributed,
  the expected log-change per odd step is
  $\mathbf{E}[\log(1/(\nu_2+1))]=-S$ with sign corrected by the
  direction of $y$.  The polylogarithm identity follows from the series
  expansion $\operatorname{Li}_s(x)=\sum_{k=1}^\infty x^k k^{-s}$
  differentiated at $s=0$.
\end{proof}

\begin{remark}
  The $\pm S$ symmetry implies that $\log y_k$ is a zero-mean random
  walk (conditionally on equal step probabilities).  This is consistent
  with the Collatz conjecture but does not prove it: the steps are
  deterministic (not independent), $E[\text{ratio}]>1$ for the
  $\operatorname{Im}(\tau)$ process itself, and the terminal condition
  $y=\frac{1}{2}$ is not equivalent to $n=1$.
\end{remark}

%% ═══════════════════════════════════════════════════════════════════════════
\section{Directions for further work}
\label{sec:further}
%% ═══════════════════════════════════════════════════════════════════════════

\begin{enumerate}[leftmargin=*, label=\textbf{\arabic*.}]
  \item \textbf{Explicit fiber construction.}
        For each $n\in\N$ and each $\ell\geq 1$, there exists $r$ with
        $3^r n \equiv 1\pmod{2^\ell}$ (by Euler's theorem).  A proof
        of the conjecture via~\ref{it:C} would need to show that, for
        such $(r,\ell)$, the target $b^* := 1-3^r n/2^\ell$ lies in
        $\mathcal{B}(r,\ell)$.  A natural approach is to show that
        $b^*$ is always expressible as an ordered partial sum of the
        form \eqref{eq:intercept}.

  \item \textbf{Minimum gap in $\mathcal{B}(r,\ell)$.}
        Quantifying $\min_{b\neq b'\in\mathcal{B}(r,\ell)}|b-b'|$
        would bound how large $(r,\ell)$ must be for the target
        $b^*$ to have any chance of lying in $\mathcal{B}(r,\ell)$.
        Since $|\mathcal{B}(r,\ell)|=\binom{r+\ell}{r}$ and all
        elements have denominator $\leq 2^\ell$, the minimum gap
        is at least $2^{-\ell}$.

  \item \textbf{The $(2,3)$-Thompson group $G_{(2,3)}$.}
        The group of piecewise-linear homeomorphisms of $[0,1]$ with
        breakpoints in $\Zs$ and slopes in $\{2^s 3^t\}$ is the
        natural rigid ambient group; it contains $\mathcal{S}$ as an
        infinite-type element.  Establishing that its derived subgroup
        is simple and computing its abelianisation (conjectured $\Z^2$,
        using the Stein--Brown framework~\cite{Stein1992}) would be
        the first new result in this direction.

  \item \textbf{C*-simplicity.}
        If $G_{(2,3)}$ is C*-simple, any action on a compact space has
        a unique stationary measure, strongly constraining the orbit
        structure.  This could close the amenability gap by ruling out
        exceptional orbits.

  \item \textbf{Furstenberg's $\times 2$, $\times 3$ connection.}
        The slope group $\{2^s 3^t\}$ is the same multiplicative
        structure appearing in Furstenberg's conjecture on
        simultaneously $\times 2$- and $\times 3$-invariant measures
        on $\R/\Z$.  Recent progress by
        Shmerkin~\cite{Shmerkin2019} and Host~\cite{Host1995} may
        be relevant to understanding the distribution of the semigroup
        action.
\end{enumerate}

%% ═══════════════════════════════════════════════════════════════════════════
\section*{Acknowledgements}
%% ═══════════════════════════════════════════════════════════════════════════

These results emerged from a computational investigation beginning with
a Godel-encoding approach (Section~\ref{sec:tau}), proceeding through
Im$(\tau)$ dynamics and the $\pm S$ symmetry, and culminating in the
affine embedding of Sections~\ref{sec:embedding}--\ref{sec:gap}.
All algebraic computations were verified exactly using rational
arithmetic over $\Zs$.

%% ═══════════════════════════════════════════════════════════════════════════
\bibliographystyle{alpha}
\begin{thebibliography}{10}

\bibitem{BelkBrown2005}
J.~Belk and K.~S. Brown.
\newblock Forest diagrams for elements of Thompson's group $F$.
\newblock \emph{Internat.\ J.\ Algebra Comput.}, 15(5--6):815--850, 2005.

\bibitem{BrunnerSidki1999}
A.~M. Brunner and S.~Sidki.
\newblock The generation of $\mathrm{GL}(n,\mathbb{Z})$ by finite state
  automata.
\newblock \emph{Internat.\ J.\ Algebra Comput.}, 9(3--4):421--463, 1999.

\bibitem{Cleary2000}
S.~Cleary.
\newblock Regular subdivision in $\mathbb{Z}[1/n]$ and the Thompson groups.
\newblock \emph{J.\ Algebra}, 225(2):714--736, 2000.

\bibitem{Day1957}
M.~M. Day.
\newblock Amenable semigroups.
\newblock \emph{Illinois J.\ Math.}, 1:509--544, 1957.

\bibitem{Haagerup2016}
U.~Haagerup and M.~Olesen.
\newblock The Thompson groups $T$ and $V$ are not inner amenable.
\newblock \emph{Pacific J.\ Math.}, 272(1):129--166, 2014.

\bibitem{Host1995}
B.~Host.
\newblock Nombres normaux, entropie, translations.
\newblock \emph{Israel J.\ Math.}, 91(1--3):419--428, 1995.

\bibitem{Lagarias2010}
J.~C. Lagarias (ed.).
\newblock \emph{The Ultimate Challenge: The $3x+1$ Problem}.
\newblock American Mathematical Society, Providence, RI, 2010.

\bibitem{Lindenstrauss2001}
E.~Lindenstrauss.
\newblock Pointwise theorems for amenable groups.
\newblock \emph{Invent.\ Math.}, 146(2):259--295, 2001.

\bibitem{Shmerkin2019}
P.~Shmerkin.
\newblock On Furstenberg's intersection conjecture, self-similar measures,
  and the $L^q$ norms of convolutions.
\newblock \emph{Ann.\ of Math.\ (2)}, 189(2):319--391, 2019.

\bibitem{Stein1992}
M.~Stein.
\newblock Groups of piecewise linear homeomorphisms.
\newblock \emph{Trans.\ Amer.\ Math.\ Soc.}, 332(2):477--514, 1992.

\bibitem{Tao2019}
T.~Tao.
\newblock Almost all orbits of the Collatz map attain almost bounded values.
\newblock \emph{Forum Math.\ Pi}, 10:e12, 2022.

\bibitem{Terras1976}
R.~Terras.
\newblock A stopping time problem on the positive integers.
\newblock \emph{Acta Arith.}, 30(3):241--252, 1976.

\bibitem{Yolcu2023}
E.~Yolcu, S.~Aaronson, and M.~J.~H. Heule.
\newblock A concrete realization of the Collatz conjecture in formal rewriting
  systems.
\newblock \emph{J.\ Automated Reasoning}, 67(2), 2023.

\end{thebibliography}

\end{document}
